\cvsection{Expériences professionnelles}

\begin{cventries}
  \cventry
    {Stagiaire ingénieur de recherche en IA}
    {INRIA}
    {Bordeaux, France}
    {Mars 2026 -- Septembre 2026}
    {
      \begin{cvitems}
        \item {Transformation d’un besoin défini avec mon tuteur en une application fonctionnelle aidant les chercheurs à trouver des articles scientifiques.}
        \item {Conception et développement en autonomie, avec le suivi de mes encadrants, d’un agent IA multi-étapes et des outils nécessaires à ses différentes actions.}
        \item {Construction de l’agent sans framework d’orchestration afin de maîtriser l’exécution des outils, l’état, la mémoire, le contexte et les enchaînements de tâches.}
        \item {Conception et ajustement de prompts pour personnaliser le comportement de l’agent et guider la sélection et l’utilisation de ses outils.}
        \item {Utilisation de Codex et de DeepSeek dans des workflows de développement assisté par l’IA pour accompagner le prototypage et l’amélioration du code.}
        \item {Développement du backend en Python/FastAPI, d’API REST documentées avec OpenAPI et de la persistance MongoDB.}
        \item {Développement d’une interface React/Next.js permettant aux chercheurs d’interagir avec l’agent et ses outils.}
        \item {Intégration de l’API OpenAI et mise en œuvre d’un pipeline RAG avec embeddings, recherche sémantique et ChromaDB.}
        \item {Création d’un outil de collecte automatisée de données et de webhooks transmettant les résultats de tâches cron.}
        \item {Conteneurisation de composants avec Docker et mise en place d’un pipeline CI/CD sur GitHub.}
        \item {Définition d’un protocole d’évaluation, analyse des résultats et contribution à un article scientifique comparant plusieurs LLM.}
      \end{cvitems}
    }

  \cventry
    {Stagiaire de recherche en apprentissage par renforcement}
    {Université Karunya}
    {Coimbatore, Tamil Nadu, Inde}
    {Juin 2025 -- Septembre 2025}
    {
      \begin{cvitems}
        \item {Développement d’un modèle d’apprentissage par renforcement et de son environnement d’entraînement pour la navigation autonome d’un robot hospitalier.}
        \item {Déploiement et optimisation du logiciel et du modèle sur un Raspberry Pi embarqué.}
      \end{cvitems}
    }
\end{cventries}
